\begin{solution}{38}
    \begin{subsolution}
        一维麦克斯韦速率分布为
        \begin{equation}
            \mathrm{dn} = n \sqrt {\frac {m}{2 \pi k _ {\mathrm{B}} T}} \mathrm{e} ^ {- \frac {m v _ {x} ^ {2}}{2 k _ {\mathrm{B}} T}} \mathrm{dv} _ {x}
            \label{eq:1.s38}
        \end{equation}
        由于另外两个方向的速度分量不影响通量，因此无需考虑另两个分量，从而
        \begin{equation}
            \varPhi = \int v _ {x} \mathrm{d} n
            \label{eq:2.s38}
        \end{equation}
        将 \eqref{eq:1.s38} 式代入得
        \begin{equation}
            \varPhi = \int_ {0} ^ {+ \infty} n \sqrt {\frac {m}{2 \pi k _ {\mathrm{B}} T}} \mathrm{e} ^ {- \frac {m v _ {x} ^ {2}}{2 k _ {\mathrm{B}} T}} v _ {x} \mathrm{d} v _ {x} = n \sqrt {\frac {m}{2 \pi k _ {\mathrm{B}} T}} \frac {1}{2 \cdot \frac {m}{2 k _ {\mathrm{B}} T}} = n \sqrt {\frac {k _ {\mathrm{B}} T}{2 \pi m}}
            \label{eq:3.s38}
        \end{equation}
        进而利用 $P = nk_{\mathrm{B}}T$ 可得
        \begin{equation}
            \bar {\phi} = \frac {P}{\sqrt {2 \pi m k _ {\mathrm{B}} T}}
            \label{eq:4.s38}
        \end{equation}
        本小问也可以用泄流公式 $\Gamma = \frac{1}{4} n\overline{v}$ 导出。
    \end{subsolution}
    \begin{subsolution}
        \begin{subsubsolution}
            实际上这只是个参考系变换的问题。地面系的速度
            \begin{equation}
                \boldsymbol {v} = \boldsymbol {v} ^ {\prime} + u \boldsymbol {e} _ {x} \implies \mathrm{d} ^ {3} \boldsymbol {v} = \mathrm{d} ^ {3} \boldsymbol {v} ^ {\prime}
                \label{eq:5.s38}
            \end{equation}
            其中 $\boldsymbol{v}'$ 为质心系中气体分子的速度。进而
            \begin{equation}
                \begin{array}{r l r} & & {f (\pmb {v}) = f _ {0} (\pmb {v ^ {\prime}}) = \left(\frac {m}{2 \pi k _ {\mathrm{B}} T}\right) ^ {3 / 2} \exp \left(- \frac {m (\pmb {v} - u \pmb {e} _ {x}) ^ {2}}{2 k _ {\mathrm{B}} T}\right)} \\ & & {\approx \left(\frac {m}{2 \pi k _ {\mathrm{B}} T}\right) ^ {3 / 2} \left(1 + \frac {m}{k _ {\mathrm{B}} T} \pmb {v} \cdot u \pmb {e} _ {x}\right) \exp \left(- \frac {m \pmb {v} ^ {2}}{2 k _ {\mathrm{B}} T}\right) = f _ {0} (\pmb {v}) \left(1 + \frac {m}{k _ {\mathrm{B}} T} v _ {x} u\right)} \end{array}
                \label{eq:6.s38}
            \end{equation}
            由此可见
            \begin{equation}
                \alpha (\pmb {v}) = \frac {m}{k _ {\mathrm{B}} T} v _ {x}
                \label{eq:7.s38}
            \end{equation}
        \end{subsubsolution}
        \begin{subsubsolution}
            取壁面切平面，法线为 $y$ 轴，撞击壁面需 $v_{y} < 0$。气体向壁面传递的轴向动量通量为
            \begin{equation}
                M _ {x} = \int_ {v _ {y} <   0} (m v _ {x}) (- v _ {y}) f (\pmb {v}) \mathrm{d} ^ {3} \pmb {v}
                \label{eq:8.s38}
            \end{equation}
            代入 $f(\pmb {v})\approx f_0(\pmb {v})\left(1 + \frac{mv_xu}{k_{\mathrm{B}}T}\right)$，其中
            \begin{equation}
                f _ {0} (\boldsymbol {v}) = n \left(\frac {m}{2 \pi k _ {\mathrm{B}} T}\right) ^ {3 / 2} \exp \left(- \frac {m (v _ {x} ^ {2} + v _ {y} ^ {2} + v _ {z} ^ {2})}{2 k _ {\mathrm{B}} T}\right)
                \label{eq:9.s38}
            \end{equation}
            并注意到 $\int_{v_y < 0} mv_x(-v_y)f_0\mathrm{d}^3\pmb {v} = 0$，得到
            \begin{equation}
                M _ {x} = \frac {m ^ {2} u}{k _ {\mathrm{B}} T} \int_ {v _ {y} <   0} v _ {x} ^ {2} (- v _ {y}) f _ {0} (\pmb {v}) \mathrm{d} ^ {3} \pmb {v}
                \label{eq:10.s38}
            \end{equation}
            令 $\alpha = \frac{m}{2k_{\mathrm{B}}T}$，三方向积分分离：
            \begin{equation}
                \int_ {- \infty} ^ {0} (- v _ {y}) \mathrm{e} ^ {- \alpha v _ {y} ^ {2}} \mathrm{d} v _ {y} = \frac {1}{2 \alpha}, \quad \int_ {- \infty} ^ {\infty} \mathrm{e} ^ {- \alpha v _ {z} ^ {2}} \mathrm{d} v _ {z} = \sqrt {\frac {\pi}{\alpha}}, \quad \int_ {- \infty} ^ {\infty} v _ {x} ^ {2} \mathrm{e} ^ {- \alpha v _ {x} ^ {2}} \mathrm{d} v _ {x} = \frac {\sqrt {\pi}}{2 \alpha^ {3 / 2}}
                \label{eq:11.s38}
            \end{equation}
            代回式 \eqref{eq:10.s38} 并与归一化因子合并，得到
            \begin{equation}
                M _ {x} = P u \sqrt {\frac {m}{2 \pi k _ {\mathrm{B}} T}}
                \label{eq:12.s38}
            \end{equation}
            完全漫反射下再发射的平均轴向动量为零，因此壁面对气体的轴向力密度为 $-M_{x}$。单位长度管段的壁面面积为 $2\pi a$，故阻力（取 $+x$ 为正）为
            \begin{equation}
                F _ {\mathrm{drag}} = - 2 \pi a M _ {x} = - a P u \sqrt {\frac {2 \pi m}{k _ {\mathrm{B}} T}}
                \label{eq:13.s38}
            \end{equation}
        \end{subsubsolution}
    \end{subsolution}
    \begin{subsolution}
        在 $u = 0$ 时, 取垂直于 $x$ 轴的截面, 由（1）中“平面法线取为 $x$”的通量公式 \eqref{eq:4.s38} 式可知零净数通量条件为
        \begin{equation}
            \frac {P (x)}{\sqrt {T (x)}} = \frac {P (x + \mathrm{d} x)}{\sqrt {T (x + \mathrm{d} x)}}
            \label{eq:14.s38}
        \end{equation}
        对 \eqref{eq:14.s38} 式取微分得
        \begin{equation}
            \frac {1}{P} \frac {\mathrm{d} P}{\mathrm{d} x} - \frac {1}{2 T} \frac {\mathrm{d} T}{\mathrm{d} x} = 0, \quad \frac {\mathrm{d} P}{\mathrm{d} x} = \frac {P}{2 T} \frac {\mathrm{d} T}{\mathrm{d} x}
            \label{eq:15.s38}
        \end{equation}
        长度为 $\mathrm{dx}$ 的气柱两端压强差在截面积 $A = \pi a^2$ 上产生轴向合力
        \begin{equation}
            F _ {\mathrm{press}} = - A \frac {\mathrm{d} P}{\mathrm{d} x} \mathrm{d} x
            \label{eq:16.s38}
        \end{equation}
        此时 $u = 0$ 使得 $F_{\mathrm{drag}} = 0$，力平衡为
        \begin{equation}
            F _ {\mathrm{press}} + F _ {\mathrm{th}} \mathrm{dx} = 0
            \label{eq:17.s38}
        \end{equation}
        故
        \begin{equation}
            F _ {\mathrm{th}} = A \frac {\mathrm{d} P}{\mathrm{d} x} = \pi a ^ {2} \frac {P}{2 T} \frac {\mathrm{d} T}{\mathrm{d} x}
            \label{eq:18.s38}
        \end{equation}
    \end{subsolution}
    \begin{subsolution}
        一般稳态下轴向力平衡为
        \begin{equation}
            F _ {\mathrm{press}} + F _ {\mathrm{th}} \mathrm{d} x + F _ {\mathrm{drag}} \mathrm{d} x = 0
            \label{eq:19.s38}
        \end{equation}
        代入 \eqref{eq:16.s38}、\eqref{eq:18.s38}、\eqref{eq:13.s38} 式并约去 $\mathrm{d}x$:
        \begin{equation}
            - \pi a ^ {2} \frac {\mathrm{d} P}{\mathrm{d} x} + \pi a ^ {2} \frac {P}{2 T} \frac {\mathrm{d} T}{\mathrm{d} x} - a P u \sqrt {\frac {2 \pi m}{k _ {\mathrm{B}} T}} = 0
            \label{eq:20.s38}
        \end{equation}
        解得
        \begin{equation}
            u = \sqrt {\frac {k _ {\mathrm{B}} T}{2 \pi m}} \pi a \left(\frac {1}{2 T} \frac {\mathrm{d} T}{\mathrm{d} x} - \frac {1}{P} \frac {\mathrm{d} P}{\mathrm{d} x}\right)
            \label{eq:21.s38}
        \end{equation}
        净分子流率为
        \begin{equation}
            \dot {N} = n A u = \frac {P}{k _ {\mathrm{B}} T} \pi a ^ {2} u
            \label{eq:22.s38}
        \end{equation}
        将 \eqref{eq:21.s38} 式代入 \eqref{eq:22.s38} 式得
        \begin{equation}
            \dot {N} = \frac {P}{k _ {\mathrm{B}} T} \pi a ^ {2} \cdot \sqrt {\frac {k _ {\mathrm{B}} T}{2 \pi m}} \pi a \left(\frac {1}{2 T} \frac {\mathrm{d} T}{\mathrm{d} x} - \frac {1}{P} \frac {\mathrm{d} P}{\mathrm{d} x}\right) = \frac {\pi^ {3 / 2} a ^ {3}}{\sqrt {2 m k _ {\mathrm{B}} T}} \left(\frac {P}{2 T} \frac {\mathrm{d} T}{\mathrm{d} x} - \frac {\mathrm{d} P}{\mathrm{d} x}\right)
            \label{eq:23.s38}
        \end{equation}
    \end{subsolution}
\end{solution}